\scanpage{35}
INSTITUT\\
DES HAUTES ÉTUDES SCIENTIFIQUES\\
BURES-SUR-YVETTE (S.-\&-O.)\\
928-48-53

Les éléments de $\operatorname{DR}(M)$ sont les éléments de $\operatorname{DR}(M)_K\simeq V\otimes_{\mathbf Q}K$ de la forme $z=x+iy$ $(x,y\in V_{\mathbf R})$ avec $x=fx$, $y=-fy$ (on exprime $\tau z=z$ i.e. $f\sigma z=z$ i.e. $\sigma z=fz$).

\deleted{\unsure{Ceci|C’est à}} Donc l'espace $\operatorname{DR}(M)$ admet une \inserted{$\mathbf Q$-structure de} façon naturelle \deleted{\unsure{de structure}} : on prend le sous-espace des $x+iy$, $x,y\in V$, $x=fx$, $y=-fy$. Soit $W$ cet espace, il mérite le nom de \uline{sous-espace de Betti} \inserted{\uline{tordu}} de $\operatorname{DR}(M)$.

Il est stable par $f$. Mais la formation de ce sous-espace \unsure{se comporte-t-il} comme il faut par $\otimes$ ? \uline{Oui}.

\diagram{\begin{tikzcd}[ampersand replacement=\&,column sep=small,row sep=large]
H^{\unsure{*|0}}(X(K),\mathbf Q)=V
 \arrow[r,phantom,"\subset" description]
 \arrow[d,phantom,"\cup" description]
\& \operatorname{DR}(M)_K\simeq\operatorname{DR}(M_K)
 \arrow[d,phantom,"\cup" description]\\
H^{\unsure{*|0}}(X(\mathbf R),\mathbf Q)=V\cap W=V^f=W^{\unsure{F|f}}
 \arrow[r,phantom,"\subset" description]
\& W
\end{tikzcd}}

On a ainsi un deuxième foncteur fibre sur $\mathcal M$ à valeurs dans $\mathbf Q$, mais il est isom \inserted{comme foncteur \unsure{$\oplus$-isant|$\otimes$-isant}} au premier, \deleted{\unsure{en effet}} comme on voit en écrivant $V\simeq V^++V^-$, $W=W^++W^-$, et utilisant les isomorphismes $V^+\simeq W^+$, $V^-\simeq W^-$ canoniques. Cet isom $V\unsure{-|\longrightarrow}W$ est \inserted{\unsure{$f$.}} induit par l'autom de $\operatorname{DR}(M_K)\simeq V_K\simeq W_K$ qui est l'identité sur $V_K^+$, et la multiplication par $i$ sur $V_K^-$ ; Cependant, cet isom \uline{n'est} \uline{pas} compatible avec le produit \uline{tensoriel}.

Il semble que $\operatorname{Betti}^T$ et $\operatorname{Betti}$ ne sont \unsure{pas $\otimes$ isom}, l'un est déduit de l'autre en tordant par le cocycle défini par $f_{\mathbf R}$ ----
\margin{Betti tordu}
\mspage{34}
